\subsection{一重零境界の全排除と共鳴例外の二曲線への還元}
\begin{proposition}
非退化正top・対称secant負top・正Pfaffianの必要atlasでは、$aq=0$ の全領域にH3対象はない。
また $q=1$ の共鳴面 $2t+\rho^2-1=0$ に残る未知は $\rho^2=-1,-3$ の二つの数体上一パラメータ曲線だけである。
\end{proposition}
\begin{proof}
まず二重零環のHilbert級数を全値で保証する。
$T=k\langle u,v\rangle/(u^2v-vu^2,uv^2-v^2u)$ とする。
$z=uv+vu$ を導入すると
\[
 T=k[u,z][v;\sigma,\delta],\quad
 \sigma(u)=-u,\quad\sigma(z)=z,\quad\delta(u)=z,\quad\delta(z)=0
\]
というOre表示になる。両表示の写像は互いに逆であり、$T$ はdomainである。
$\deg u=\deg v=1,\deg z=2$ から
\[
 H_T(h)=\frac1{(1-h)^2(1-h^2)}.
\]
全対称負二次式は $T$ 内で中心的であり、負二平面の全GL作用はこの表示の代数に自己同型として作用する。

二重零環の逆作用 $F,F^{-1}$ と $A_2,B_2=F(A_2)$ は前節のものとする。
$\sigma_0=F^{-1}$、$a_0=F(A_2)$ として、skew Laurent環
$T[X,X^{-1};\sigma_0]$ 内で
\[
 x=X,\qquad y=a_0X^{-1}
\]
と置けば $xy=A_2,yx=B_2$ となる。
$A_2$ は非零中心元であり、$T$ はdomainなので、標準形
\[
 T\ \oplus\!\bigoplus_{n\ge1}(Tx^n\oplus Ty^n)
\]
は一次独立である。これは一般化Weyl代数のPBW表示そのものである。
元の六二次関係から $xy,yx$ と負一次式の可換性、従って $T$ の二つの三次関係が従う。
逆にこのWeyl表示は元の六関係を満たすため、二重零環と同型である。
従って二重零環は全値でdomain、Hilbert級数は
\[
 H_T(h)\frac{1+h}{1-h}=\frac1{(1-h)^4}
\]
となる。AS4十分性を追加で仮定する必要はない。

次に一重零を二重零へフィルター付きで比較する。
$a=0,q\ne0$ とし、正torusで $q=1$ とする。
$\deg_{\rm aux}(x,y,u,v)=(1,-1,0,0)$ による増大filtrationを使う。
各通常次数ではこのfiltrationは有限であり、乗法に適合する。
$q$ を含む二つの奇関係では、その項の補助次数は残る項より2低い。
従って元六関係の先頭関係は正確に $a=q=0$ の六関係となる。
二重零環から $\operatorname{gr}B$ へ全射があり、$B$ がAS4候補なら両者の通常Hilbert級数が一致するので同型である。

正常条件を補助次数の先頭へ送ると、正常核の随伴次数付き空間は二重零環の正常核へ単射となる。
前節の補助次数解析は、二重零環のdomain性と負二次式の独立性だけで成立する。
従って $\operatorname{tr}F\ne1$ なら一重零環の正常核も零、$\operatorname{tr}F=1$ なら次元は高々2である。
共鳴時には前命題の全値因子pencilが次元2を与えるため、これが正常核全体となり全切断が分解する。
$\operatorname{tr}F=2t/(1-\rho^2)$ なので共鳴式も一致する。
$q=0,a\ne0$ は座標交換でこの場合へ移り、二重零自体も既に排除済みである。
これで $aq=0$ の全値を閉じた。
特に共鳴面では $a=-\alpha D/4$ なので $D=0$ は排除済みである。

最後に残るminorの零点を全て調べる。
前命題の正常8minorに三つを追加する。
$\rho(\rho^2-1)(\rho^2+1)(\rho^2+3)D\ne0$ で使用汎関数は全て正則であり、
三minorは単元因子を除いて
\[
 (\rho+3)E,\qquad s^2F_4,\qquad s^2J_6
\]
となる。ここで
\[
 F_4=(\rho+1)^2(\rho^2+2\rho+3)+6s^2,\qquad
 J_6=(\rho+1)^2(\rho^2+1)^2+4\rho^2s^2.
\]
$E,F_4,J_6$ のidealは $\bigl(s^2,(\rho+1)^2\bigr)$ と一致する。
従って $s\ne0,\rho\ne-3$ に共通零点はない。
$\rho=-3$ でも $F_4=6s^2+24$ と $J_6=36s^2+400$ は互いに素である。
$s=0$ では $E=(\rho+1)^4(\rho^2+1)$ なので唯一の残余は $(\rho,s)=(-3,0)$。
この一点も元の全関係を再消去すると、四次商次元35、正自由三次の正常写像rank8である。
全値因子pencilと合わせて排除できる。
これで前命題の余分な $\rho^2-9$ と $E$ を除去し、$D=0$ も一重零の議論で閉じた。

\texttt{verify\_cubic\_order2\_22\_secant\_resonance\_cover.py}
は先頭六関係、Hilbert恒等式、$a=-\alpha D/4$、三minor、汎関数の極、
三多項式のidealと唯一の点を再生する。
Ore及びWeyl表示による構造証明を、有限数のrank標本だけで置き換えるものではない。
\end{proof}
